% !TEX program = xelatex
% Identity as Accountability Chain — Paraxiom Technologies
% Build: latexmk -xelatex main.tex   (or: xelatex main.tex twice)

\documentclass[11pt,a4paper,parskip=half-,DIV=12,headings=standardclasses]{scrartcl}

% =============================================================================
% Fonts (xelatex + system fonts on macOS)
% =============================================================================
\usepackage{fontspec}
\usepackage{unicode-math}
\setmainfont{Charter}[Ligatures=TeX]
\setsansfont{Helvetica Neue}[
  Scale=MatchLowercase,
  UprightFont=* Light,
  BoldFont=* Medium,
  ItalicFont=* Light Italic,
  BoldItalicFont=* Medium Italic]
\setmonofont{Menlo}[Scale=MatchLowercase]
\setmathfont{Latin Modern Math}

% =============================================================================
% Layout & typography
% =============================================================================
\usepackage{microtype}
\usepackage[margin=2.5cm,top=2.8cm,bottom=2.8cm]{geometry}
\linespread{1.15}

\addtokomafont{section}{\sffamily\bfseries\large}
\addtokomafont{subsection}{\sffamily\bfseries\normalsize}
\addtokomafont{subsubsection}{\sffamily\normalsize\itshape}
\addtokomafont{descriptionlabel}{\sffamily\bfseries}
\addtokomafont{caption}{\small}
\addtokomafont{captionlabel}{\sffamily\bfseries\small}

% =============================================================================
% Colors, tables, code, figures
% =============================================================================
\usepackage[table]{xcolor}
\definecolor{ncgray}{HTML}{505050}
\definecolor{nclight}{HTML}{F4F4F4}
\definecolor{ncaccent}{HTML}{2C5282}
\definecolor{ncrule}{HTML}{B8B8B8}
\definecolor{nckeyword}{HTML}{0F62FE}
\definecolor{nccomment}{HTML}{6F7C8A}
\definecolor{ncstring}{HTML}{8A4F2A}
\definecolor{ncwarn}{HTML}{B7410E}

\usepackage{booktabs}
\usepackage{array}
\usepackage{tabularx}
\usepackage{longtable}
\renewcommand{\arraystretch}{1.15}

\usepackage{listings}
\lstset{
  basicstyle=\ttfamily\footnotesize,
  keywordstyle=\color{nckeyword},
  commentstyle=\color{nccomment}\itshape,
  stringstyle=\color{ncstring},
  showstringspaces=false,
  breaklines=true,
  frame=single,
  framerule=0pt,
  rulecolor=\color{ncrule},
  backgroundcolor=\color{nclight},
  xleftmargin=0.5em,
  xrightmargin=0.5em,
  numbers=none,
  captionpos=b,
}
\lstdefinelanguage{rust}{
  morekeywords={fn,let,mut,struct,enum,impl,trait,pub,use,mod,where,Self,self,
                async,await,move,box,dyn,as,if,else,match,for,while,loop,return,
                type,const,static,ref,in,Option,Result,Vec,String,bool,u8,u32,u64,i32,i64,
                BoundedVec,H256,BlockNumber,AccountId,Falcon512,Sphincs,Blake2},
  morecomment=[l]{//},
  morecomment=[s]{/*}{*/},
  morestring=[b]"
}

% =============================================================================
% Math, theorems, proofs
% =============================================================================
\usepackage{amsmath}
\usepackage{amssymb}
\usepackage{amsthm}
\usepackage{mathtools}

\newtheoremstyle{ncthm}%
  {1.2ex}{1.2ex}%
  {\itshape}{}%
  {\sffamily\bfseries}{.}%
  {0.5em}{}

\newtheoremstyle{ncdef}%
  {1.2ex}{1.2ex}%
  {\normalfont}{}%
  {\sffamily\bfseries}{.}%
  {0.5em}{}

\theoremstyle{ncthm}
\newtheorem{theorem}{Theorem}[section]
\newtheorem{lemma}[theorem]{Lemma}
\newtheorem{proposition}[theorem]{Proposition}
\newtheorem{corollary}[theorem]{Corollary}

\theoremstyle{ncdef}
\newtheorem{definition}[theorem]{Definition}
\newtheorem{axiom}[theorem]{Axiom}
\newtheorem{remark}[theorem]{Remark}
\newtheorem{example}[theorem]{Example}

\renewcommand{\qedsymbol}{$\blacksquare$}

% =============================================================================
% TikZ for diagrams
% =============================================================================
\usepackage{tikz}
\usetikzlibrary{arrows.meta,positioning,shapes.geometric,fit,backgrounds,calc}

% =============================================================================
% Hyperref + cleveref (last)
% =============================================================================
\usepackage[colorlinks=true,
            linkcolor=ncaccent,
            urlcolor=ncaccent,
            citecolor=ncaccent]{hyperref}
\usepackage[capitalize]{cleveref}

% =============================================================================
% Custom commands
% =============================================================================
\newcommand{\identity}{\mathcal{I}}
\newcommand{\substrate}{\mathcal{S}}
\newcommand{\founding}{\mathcal{F}}
\newcommand{\accretion}{\mathcal{A}}
\newcommand{\relations}{\mathcal{R}}
\newcommand{\layers}{\mathcal{L}}
\newcommand{\pqset}{\Sigma_{\mathrm{PQ}}}
\newcommand{\Negl}{\mathrm{negl}}
\newcommand{\KDF}{\mathsf{KDF}}
\newcommand{\Hash}{\mathsf{H}}
\newcommand{\Sign}{\mathsf{Sign}}
\newcommand{\Verify}{\mathsf{Verify}}
\newcommand{\Enc}{\mathsf{Enc}}
\newcommand{\Dec}{\mathsf{Dec}}
\newcommand{\IPFS}{\mathsf{IPFS}}
\newcommand{\Chain}{\mathsf{Chain}}
\newcommand{\paraxiom}{\textsc{Paraxiom}}
\newcommand{\qh}{\textsc{QuantumHarmony}}
\newcommand{\nc}{\textsc{NeoCarbone}}

% =============================================================================
% Title block
% =============================================================================
\title{
  \sffamily\bfseries
  Identity as Accountability Chain\\[0.4ex]
  \large\mdseries A Post-Quantum Architecture for Persistent Identity\\
  Across Human, Machine, and AI Agent Substrates
}
\author{
  \normalsize Sylvain Cormier\\
  \footnotesize Paraxiom Technologies Inc.\\
  \footnotesize \texttt{sylvain@paraxiom.org}
}
\date{\small May 15, 2026}

% =============================================================================
\begin{document}
% =============================================================================

\maketitle

\begin{abstract}
\noindent
We propose a formal model in which \emph{identity} is not a primitive
but a layered accountability chain that accretes from a founding
individuation event and persists as long as cryptographic, relational,
and civic verification can be reconstructed. The model unifies
biological, machine, and artificial-intelligence cases under a single
schema with six accretion layers and a four-tier permanence hierarchy
for substrate-level identifiers.

We prove four results that bear directly on engineering: (i)~Founding
uniqueness --- every identity locus has exactly one origin event;
(ii)~Tier monotonicity --- the permanence of an identity claim is
non-increasing in the tier of its substrate identifier;
(iii)~Cryptographic survival --- only post-quantum-signed identities
remain verifiable past a cryptographically-relevant quantum computer
(CRQC) event; and (iv)~Crypto-shredding completeness --- master-key
deletion renders all envelope-encrypted credentials computationally
inaccessible without recourse to unpinning.

We survey the state of digital identity as of May 2026, noting that no
major civic identity system runs post-quantum signatures in production
despite regulatory deadlines of 2030--2035, and that AI agent identity
is a standards vacuum.

We then present the \paraxiom\ architecture: a three-tier system
combining on-chain post-quantum primitives (\qh, with SPHINCS+ and
Falcon-512), IPFS for bulk encrypted payloads, and Postgres+Redis
mirrors for query-layer access. The architecture is grounded in five
identity strata (S1--S5) covering humans, scoped delegations, vault
credentials, agent instances, and cross-tenant aggregates, with
provision for future extensions to swarm and protocol identities.

The treatise concludes that identity is engineering, not philosophy:
the accountability chain does not extend itself; the post-quantum
signatures do not sign themselves. What we build now determines
whether identity remains a coherent concept past 2040.
\end{abstract}

\bigskip
\noindent\textbf{Keywords:}\;
digital identity, post-quantum cryptography, accountability,
agent identity, substrate attestation, decentralized identity,
SPHINCS+, Falcon-512, Substrate, IPFS, eIDAS, AAMVA, MCP,
verifiable credentials.

\bigskip

\tableofcontents

\bigskip

% =============================================================================
\section{Introduction}
% =============================================================================

\subsection{Premise}

A governance application requires an ``identity primitive.'' A
blockchain pallet needs a stratum field. A credential vault needs a
user. Before we choose data structures, we must answer a prior
question: what \emph{is} identity, and when does it begin?

This treatise advances a single claim:

\begin{quote}
\itshape
Identity is not a primitive. It is a layered accountability chain
that accretes from the moment a locus of action becomes individuated,
and persists exactly as long as the chain can verify continuity.
When the chain breaks, identity dissolves --- regardless of whether
the substrate (body, hardware, weights file) is still running.
\end{quote}

This claim holds for humans, for processes, for AI agents, for swarms,
for organizations. The substrate differs; the accretion logic is
identical. We will state this formally
(\Cref{def:identity-locus,def:accretion-chain}), prove that it
implies four engineering invariants
(\Cref{thm:founding-uniqueness,thm:tier-monotonicity,thm:cryptographic-survival,thm:crypto-shredding}),
and demonstrate that the \paraxiom\ architecture realizes them.

\subsection{Why now}

Three forces drive this work to operational urgency in 2026.

\textbf{Post-quantum cryptography becomes load-bearing.} NIST published
ML-KEM, ML-DSA, and SLH-DSA as FIPS standards on 13~August
2024~\cite{nist-fips-203,nist-fips-204,nist-fips-205}. The U.S.\ NSA's
CNSA~2.0 prohibits classical algorithms in National Security Systems
by 2035. The EU joint PQC roadmap (June 2025) targets identity-wallet
PKI migration by the same horizon. Yet, as of May 2026, no major civic
identity system operationally runs post-quantum signatures in
production. The window in which a classical signature today can be
forged retroactively by a sufficient quantum computer is the window
in which the entire pre-CRQC accountability chain collapses.

\textbf{AI agent identity is an open standards vacuum.} The Linux
Foundation's Agentic AI Foundation launched in December~2025 with
contributions from Anthropic (MCP), OpenAI (AGENTS.md), and Block
(Goose). MCP added a server-identity primitive in
November~2025~\cite{mcp-spec-2025-11}. Google's A2A protocol uses
workload-identity federation~\cite{google-a2a}. Cloudflare-backed
IETF Web Bot Auth proposes signed HTTP requests. No standard exists
for model substrate attestation: cryptographically proving that a
given output originated from a specific weights binary on a specific
hardware root. The field is where SSI was in 2018.

\textbf{Civic identity is fragmenting.} EU eIDAS~2.0 mandates wallet
deployment by 24~December~2026~\cite{eidas-2024-1183}. US mDL
proliferates state by state with no federal coordination. The UK
GOV.UK One~Login operates uncertified against its own trust
framework. Canada lacks a federal equivalent to EUDI. The conditions
for an architecture-of-record are present.

\subsection{Method}

\Cref{sec:formal-model,sec:layer-theory,sec:founding,sec:machine-storage,sec:agent,sec:llm,sec:crypto-substrate}
develop the formal model. \Cref{sec:landscape} surveys the deployed
state of digital identity in May~2026. \Cref{sec:paraxiom-architecture}
presents the engineering realization, with proofs of crypto-shredding
completeness and verification-time accountability enforcement.
\Cref{sec:future,sec:roadmap,sec:conclusion} discuss scenario
projections, implementation sequencing, and concluding remarks.
Appendices give detailed proofs, pallet API specifications, and
sources.

% =============================================================================
\section{Formal model of identity}\label{sec:formal-model}
% =============================================================================

We begin by stating the model abstractly. Subsequent sections
specialize to biological, machine, agent, and LLM cases.

\subsection{Primitive sets}

Fix the following primitive sets:
\begin{itemize}
\item $\mathbb{T}$ --- time, totally ordered.
\item $\mathbb{K}$ --- cryptographic keypairs over some signature
  scheme $\Sigma$, with operations
  $\Sign_{sk}, \Verify_{pk}$ satisfying EUF-CMA security.
\item $\mathbb{C}$ --- claim content (arbitrary byte-strings).
\item $\mathbb{S}$ --- substrates: a typed union of biological organisms,
  physical hosts, container images, weights blobs, and similar
  individuated material configurations.
\item $\mathbb{L} = \{L_0, L_1, L_2, L_3, L_4, L_5\}$ --- the six
  accretion layers defined in \Cref{sec:layer-theory}.
\end{itemize}

\subsection{Definitions}

\begin{definition}[Founding event]\label{def:founding-event}
A \emph{founding event} is a triple
$\founding = (s, t_0, \pi)$
where $s \in \mathbb{S}$ is a substrate, $t_0 \in \mathbb{T}$ is the
event timestamp, and $\pi$ is a finite set of witnesses (the empty
set is permitted but reduces verifiability). The founding event marks
the transition of $s$ from a non-individuated state to an
individuated locus capable of accreting claims.
\end{definition}

\begin{definition}[Accretion chain]\label{def:accretion-chain}
An \emph{accretion chain} relative to founding event
$\founding = (s, t_0, \pi)$ is a finite sequence
\[
\accretion = (a_1, a_2, \ldots, a_n),
\quad
a_i = (L_i, t_i, c_i, \sigma_i, \rho_i)
\]
where $L_i \in \mathbb{L}$ is the layer tag, $t_i \geq t_0$ is the
claim timestamp, $c_i \in \mathbb{C}$ is the claim content,
$\sigma_i$ is either a signature in $\mathbb{K}$ or a relational
witness, and $\rho_i \in \{1,\ldots,i-1\} \cup \{\bot\}$ is an
optional reference to a prior claim providing context. The sequence
is monotone in time: $i < j \Rightarrow t_i \leq t_j$.
\end{definition}

\begin{definition}[Identity locus]\label{def:identity-locus}
An \emph{identity locus} is a quadruple
\[
\identity = (\substrate, \founding, \accretion, \relations)
\]
where $\substrate \in \mathbb{S}$ is its substrate,
$\founding$ is its founding event,
$\accretion$ is its accretion chain, and
$\relations \subseteq \mathbb{I} \times \mathbb{I}$ is its
relational graph with other identity loci
(where $\mathbb{I}$ is the universe of all identity loci).
\end{definition}

\begin{definition}[Verifiable identity]\label{def:verifiable-identity}
An identity locus $\identity$ is \emph{verifiable at time $t$} iff,
for every layer-3 (cryptographic) claim $a_i \in \accretion$ with
$t_i \leq t$, there exists a key $\KDF(\identity, t_i) \in \mathbb{K}$
such that $\Verify(\KDF(\identity, t_i), c_i, \sigma_i) = 1$ and the
substrate $\substrate$ retains continuity through the interval
$[t_0, t]$.
\end{definition}

\subsection{Axioms}

We posit three axioms that scope the model.

\begin{axiom}[Founding precedence]\label{ax:founding-precedence}
For every identity locus $\identity$, every claim
$a_i \in \accretion$ satisfies $t_i \geq t_0$, where
$t_0$ is the founding timestamp.
\end{axiom}

\begin{axiom}[Substrate continuity or explicit succession]\label{ax:substrate-continuity}
An identity locus persists from time $t_a$ to $t_b > t_a$ iff either
(i)~its substrate $\substrate$ remains individuated and physically
continuous over $[t_a, t_b]$, or (ii)~there exists a successor
mapping $\Phi: \identity \to \identity'$ signed by an authority $A$
appearing in $\relations$, recorded at some
$t \in [t_a, t_b]$.
\end{axiom}

\begin{axiom}[Layer independence]\label{ax:layer-independence}
The six layers in $\mathbb{L}$ are pairwise distinguishable: no claim
$a_i$ in an accretion chain may simultaneously have two layer tags.
The failure modes of different layers are statistically independent
absent explicit coupling (e.g., a single physical device holding both
Layer~3 and Layer~4 artifacts).
\end{axiom}

\Cref{ax:founding-precedence,ax:substrate-continuity,ax:layer-independence}
imply that identity is a well-defined object: it has a beginning,
a domain of persistence, and a coordinate system in
$\mathbb{L}$.

\subsection{First theorem: founding uniqueness}

\begin{theorem}[Founding uniqueness]\label{thm:founding-uniqueness}
For every identity locus $\identity$, exactly one founding event
$\founding$ exists.
\end{theorem}

\begin{proof}
By construction (\Cref{def:identity-locus}), $\identity$ specifies a
single $\founding$. Suppose for contradiction that two distinct events
$\founding_1, \founding_2$ both qualify. By
\Cref{ax:founding-precedence} all claims in $\accretion$ must precede
neither. If $t_1 < t_2$, then $\founding_2$ would violate the founding
precedence axiom with respect to claims indexed near $t_1$; if
$t_1 = t_2$, both events name the same individuation, hence
$\founding_1 = \founding_2$. Contradiction.
\end{proof}

\begin{remark}
\Cref{thm:founding-uniqueness} is the formal restatement of the
engineering principle that an identity in \texttt{identity-registry}
\cite{paraxiom-identity-registry} must have a uniquely-recorded
\texttt{created\_at\_block} and \texttt{created\_by\_tx\_hash}.
An identity record without a verifiable origin transaction is
malformed.
\end{remark}

% =============================================================================
\section{Layer theory}\label{sec:layer-theory}
% =============================================================================

We name and motivate the six accretion layers, then prove their
independence and ordering properties.

\subsection{The six layers}

\begin{description}
\item[$L_0$: Substrate.] The material continuity --- body, hardware,
  weights blob. Necessary precondition; identity-bearing in the sense
  that its destruction ends the locus (subject to
  \Cref{ax:substrate-continuity}).
\item[$L_1$: Narrative.] The locus's own claim of continuity:
  ``I am the same entity that did $X$ yesterday.''
  Implemented variously as autobiographical memory, persistent state on
  disk, context window, institutional documentation.
\item[$L_2$: Relational.] Other accountability loci attest that
  $\identity$ has continuity. Family, peers, federation members,
  protocol counterparties.
\item[$L_3$: Cryptographic.] A keypair (PQ-secured under the model
  assumptions of \Cref{sec:crypto-substrate}). The first layer that
  scales without prior relationship.
\item[$L_4$: Civic.] A jurisdictional attestation: birth
  certificate, passport, KYC bundle, corporate registration.
\item[$L_5$: Behavioral.] The cumulative track record of actions
  taken --- votes, decisions, signed attestations, transferred
  resources, resolved disputes. Aggregates of $L_3$ and $L_2$ over
  time.
\end{description}

\subsection{Diagram of accretion}

\begin{figure}[h]
\centering
\begin{tikzpicture}[every node/.style={font=\small},
                    layer/.style={draw=ncrule,fill=nclight,
                                  rounded corners=2pt,
                                  minimum width=10cm,
                                  minimum height=0.6cm,
                                  inner sep=4pt}]
  \node[layer] (L0) at (0,0)   {$L_0$ \;\textsf{Substrate}\;(body, hardware, weights)};
  \node[layer] (L1) at (0,0.9) {$L_1$ \;\textsf{Narrative}\;(self-attested continuity)};
  \node[layer] (L2) at (0,1.8) {$L_2$ \;\textsf{Relational}\;(witness graph)};
  \node[layer] (L3) at (0,2.7) {$L_3$ \;\textsf{Cryptographic}\;(PQ keypair)};
  \node[layer] (L4) at (0,3.6) {$L_4$ \;\textsf{Civic}\;(KYC bundle hash, jurisdictional cert)};
  \node[layer] (L5) at (0,4.5) {$L_5$ \;\textsf{Behavioral}\;(reputation tokens, activity log)};
  \draw[dashed,ncrule] (-5.5,-0.45) -- (5.5,-0.45);
  \node at (0,-0.75) {\itshape Founding individuation event};
\end{tikzpicture}
\caption{The six accretion layers above the founding individuation
event. Each layer can fail independently; combined failure of $k \geq
4$ layers causes identity dissolution.}
\label{fig:accretion}
\end{figure}

\subsection{Layer composition is non-commutative in time}

\begin{proposition}[Layer accretion order is not strict]\label{prop:layer-order}
The temporal order in which layers accrete is not fixed. An identity
may acquire $L_3$ before $L_4$ (a private key before legal
recognition), $L_2$ before $L_1$ (others recognize an organism before
the organism develops self-narrative), or $L_5$ before $L_3$ (a
reputation accumulated under pseudonym, later bound to a key).
\end{proposition}

\Cref{prop:layer-order} is consequential for design: the identity
registry must permit any layer to be backfilled, not enforce a
canonical ordering.

% =============================================================================
\section{Founding events across substrates}\label{sec:founding}
% =============================================================================

We specialize the founding-event definition to the four canonical
substrate types.

\subsection{Biological founding}

A human identity begins not at conception (continuous chemical
process) nor at issuance of the birth certificate (administrative
echo), but at \emph{first individuation}: the moment after birth at
which the organism's state is distinguishable from its parents'.
Witnesses (mother, attending professional) populate the initial
relational graph $\relations$.

\subsection{Machine founding}

A host's identity begins at first boot, when the kernel writes
\texttt{/etc/machine-id} via \texttt{systemd-machine-id-setup}, when
entropy is harvested, and when initial cryptographic material (SSH
host keys, TPM-derived keys) is established. We discuss the precise
storage locations and their permanence tiers in \Cref{sec:machine-storage}.

\subsection{Container founding}

A container's identity begins at \texttt{docker run} (or equivalent
runc/containerd invocation). The image hash supplies $L_0$
cryptographically; the container ID supplies $L_1$ as the
self-narrative anchor.

\subsection{LLM founding}

An LLM substrate's identity begins at \emph{weights initialization}:
two checkpoints trained from different random seeds on identical
data are distinct substrate identities, even when statistically
indistinguishable in behavior. An LLM \emph{instance} identity, by
contrast, begins at \emph{first inference call} on a deployed
endpoint, when the deployment ID is assigned. We elaborate the
substrate--instance distinction in \Cref{sec:llm}.

\subsection{Unifying schema}

\begin{table}[h]
\centering
\begin{tabular}{@{}lllll@{}}
\toprule
\textsf{Domain} & \textsf{Substrate $\substrate$} & \textsf{Founding $\founding$} & \textsf{First $L_1$/$L_2$ accretion}\\
\midrule
Human    & Body            & Birth                  & Mother, witnesses\\
Machine  & Hardware        & First boot             & Hostname, SSH host keys\\
Container & Image           & \texttt{docker run}    & Container ID, mount state\\
Process  & Code binary     & First invocation       & PID, parent, env\\
LLM substrate & Architecture + dataset & Weights init    & Checkpoint hash, training log\\
LLM instance  & Weights         & First inference   & Deployment ID, context\\
\bottomrule
\end{tabular}
\caption{Founding events by substrate domain.}
\label{tab:founding-schema}
\end{table}

% =============================================================================
\section{Machine identity storage}\label{sec:machine-storage}
% =============================================================================

The model is concrete: machine identity lives in named files,
registers, and hardware tokens at known locations. Understanding the
\emph{permanence tier} of each storage location is the engineering
prerequisite for trusting the resulting identity claim.

\subsection{The four-tier hierarchy}

\begin{definition}[Permanence tier]\label{def:tier}
A storage location for an identity artifact has \emph{permanence tier}
$T \in \{0,1,2,3\}$ defined as follows. Tier~$0$ identifiers are
\emph{manufacturer-fused}: present at production and unchangeable
without replacing hardware. Tier~$1$ identifiers are
\emph{installation-rooted}: established at first boot or first
install, persistent across restarts, lost on reinstall. Tier~$2$
identifiers are \emph{operationally configured}: persistent across
restarts but routinely changeable by administrators. Tier~$3$
identifiers are \emph{runtime-ephemeral}: regenerated each session.
\end{definition}

\subsection{Tier~0 (manufacturer-fused)}

\begin{itemize}
\item \textbf{SMBIOS/DMI UUID}: burned at manufacture, queryable via
  \texttt{dmidecode -s system-uuid}. Resides in firmware tables, not
  on disk.
\item \textbf{TPM 2.0 Endorsement Key (EK)}: RSA or ECC keypair burned
  into the TPM at manufacture, certified by the manufacturer's CA.
  Root of trust for hardware attestation.
\item \textbf{Apple Secure Enclave UID, Google Titan-M2 keys,
  Microsoft Pluton root keys}: per-chip UIDs fused at manufacture,
  participating in cryptographic operations without exfiltration.
\item \textbf{Storage device serials}: ATA/NVMe serials per
  physical disk, surviving reformat.
\item \textbf{Network MAC addresses}: set per NIC at manufacture but
  software-spoofable, hence straddling Tier~0 and Tier~3.
\end{itemize}

\subsection{Tier~1 (installation-rooted)}

\begin{itemize}
\item \textbf{\texttt{/etc/machine-id}}: 128-bit identifier generated
  by \texttt{systemd-machine-id-setup} on first boot. The
  \texttt{systemd} documentation explicitly states that this ID
  ``should be considered confidential and must not be exposed in
  untrusted environments,'' and that downstream applications must
  apply \texttt{sd\_id128\_get\_machine\_app\_specific()} to derive
  per-application identifiers rather than using the raw value.
\item \textbf{\texttt{/etc/ssh/ssh\_host\_\{rsa,ecdsa,ed25519\}\_key}}:
  generated at first SSH daemon configuration. These are the host's
  first cryptographic identity artifact, anchoring SSH
  trust-on-first-use and SSHFP DNS records.
\item \textbf{TPM-owned derived keys}: Storage Root Key, Attestation
  Identity Keys, platform-bound passkey credentials --- generated at
  enrollment but cryptographically descended from the Tier-0 EK.
\item \textbf{Cloud-init instance ID}: at \texttt{/var/lib/cloud/data/
  instance-id}, mirroring the cloud provider's instance handle.
\end{itemize}

\subsection{Tier~2 (operationally configured)}

\begin{itemize}
\item \textbf{\texttt{/etc/hostname}}: human-readable. Conventional
  primary identifier but cryptographically meaningless on its own.
\item \textbf{Kerberos host principal keys}
  (\texttt{/etc/krb5.keytab}); TLS server certificates; WireGuard
  keypairs; cloud IAM role bindings.
\end{itemize}

\subsection{Tier~3 (runtime-ephemeral)}

\begin{itemize}
\item PID, container ID, network state, in-kernel
  \texttt{struct cred}, namespaces, SELinux/AppArmor labels.
\end{itemize}

\subsection{Tier monotonicity theorem}

\begin{theorem}[Tier monotonicity]\label{thm:tier-monotonicity}
Let $P(T)$ denote the unconditional probability that an identifier
of permanence tier $T$ survives the time interval $[t_0, t_0 + \Delta]$
for $\Delta$ on the order of a host's installation lifetime
($\sim$1--10~years). Then
\[
P(0) \geq P(1) \geq P(2) \geq P(3),
\]
with strict inequality $P(0) > P(3)$.
\end{theorem}

\begin{proof}
By \Cref{def:tier} and operational invariants. Tier~0 identifiers
are unchangeable without hardware replacement; the probability of
hardware replacement over the operational lifetime defines $1 - P(0)$.
Tier~1 identifiers are lost on reinstall, with reinstall rate strictly
greater than hardware-replacement rate; thus $P(1) \leq P(0)$.
Tier~2 identifiers are changed administratively at rates strictly
greater than reinstall rate (hostnames, certificates, IAM bindings
rotate routinely); $P(2) \leq P(1)$. Tier~3 identifiers are lost on
reboot or session end, occurring at rates dramatically greater than
administrative configuration rotation; $P(3) \leq P(2)$ with strict
inequality typical. Strict inequality $P(0) > P(3)$ follows from the
typical many-orders-of-magnitude separation between hardware lifetime
and process lifetime.
\end{proof}

\begin{corollary}[Substrate identifier selection]\label{cor:substrate-selection}
Given a choice of substrate identifiers available at registration
time, the post-quantum identity stack should select the
\emph{lowest-numbered tier available} for the substrate identity
field. For an agent running on bare metal with TPM, this is the
TPM-quoted EK identifier; on virtualized hardware, the SMBIOS UUID;
on containers without TPM passthrough, the container image hash plus
container ID.
\end{corollary}

\Cref{cor:substrate-selection} directly informs the
\texttt{instance\_anchor} field of the agent-identity record:
\[
\texttt{instance\_anchor} := \Hash_{\textsf{Blake2-256}}(\text{tier}_{\min} \,\|\, \text{agent\_kind} \,\|\, t_{\mathrm{spawn}}),
\]
mixing the lowest-tier available identifier with runtime context to
produce a stable yet instance-specific anchor.

\subsection{Platform-specific tier locations}

\begin{table}[h]
\centering
\small
\begin{tabular}{@{}lp{4cm}p{4cm}p{4cm}@{}}
\toprule
\textsf{Tier} & \textsf{Linux} & \textsf{macOS} & \textsf{Windows}\\
\midrule
0 & SMBIOS UUID; TPM EK; SoC UID & \texttt{IOPlatformUUID}; Secure Enclave UID & SMBIOS UUID; Pluton root\\
1 & \texttt{/etc/machine-id}; SSH host keys; TPM AIKs &
    System UUID; Recovery key; Keychain enrollment &
    \texttt{HKLM\textbackslash...\textbackslash MachineGuid};
    TPM AIKs\\
2 & Hostname; krb5 keytab; TLS certs; WireGuard keys &
    Hostname; user keychain certs; configuration profiles &
    Computer SID; domain bindings\\
3 & PID; container ID; namespaces & PID; sandbox ID & PID; session token\\
\bottomrule
\end{tabular}
\caption{Permanence tiers across major operating systems.}
\label{tab:tier-platforms}
\end{table}

% =============================================================================
\section{Agent identity}\label{sec:agent}
% =============================================================================

\subsection{Three forms of agency}

A \emph{process} is the weakest agency: a running execution context
with PID, parent, and loaded code. By default it carries no $L_3$
cryptographic claim; it inherits trust from its parent.

A \emph{container} adds substrate cryptographic anchor by way of
the image hash, which itself is signed by its publisher.

A \emph{true agent}, in our sense, is a program acting on behalf of a
principal. It requires:
\begin{enumerate}
\item Its own cryptographic identity ($L_3$): a Falcon-512 or
  SPHINCS+ keypair, generated at first registration.
\item A delegation chain from the principal ($L_2$): a parent identity
  in $\relations$ tracing to an $L_4$-anchored or
  reputation-bonded $L_3$ identity.
\item Scope bounds: the set of capabilities the agent is authorized
  for.
\item Reputation ($L_5$): an accumulated track record.
\item Revocability: the principal must be able to revoke without
  cooperation from the agent.
\end{enumerate}

\subsection{The stratum lattice}

We define a partial order $\prec$ on identity loci. Write
$\identity_1 \prec \identity_2$ iff $\identity_1$ appears as a
parent identity for $\identity_2$ in $\relations$.

\begin{definition}[Stratification]\label{def:stratification}
The \emph{stratification} of an identity space partitions loci into
five strata:
\begin{description}
\item[S1 Person] KYC-rooted human identity. Root of the $\prec$ order.
\item[S2 Scoped] Person $\times$ tenant $\times$ role. Parent: an S1.
\item[S3 Vault] Encrypted credential commitments. Held by an S1, S2,
  or organization S5; not itself an identity-bearing entity but a
  record attached to one.
\item[S4 Agent instance] Running agent. Parent: any of S1, S2, or S5.
\item[S5 Cross-tenant] Organizational aggregate identity. Parent: a
  multi-signature of S1 founders or jurisdictional authority.
\end{description}
\end{definition}

\subsection{Stratum well-foundedness}

\begin{theorem}[Stratum well-foundedness]\label{thm:stratum-wellfounded}
The relation $\prec$ restricted to S2, S4, S5 identities induces a
well-founded ordering: every non-S1 identity has an S1 ancestor in
finitely many $\prec$-steps.
\end{theorem}

\begin{proof}
Assume the chain registration policy of \Cref{sec:paraxiom-architecture}:
every S2/S4 identity registration extrinsic requires a non-null
\texttt{parent\_identity} pointer, and the registration is
chain-validated against the existence of the parent. Every S5
identity registration requires either a multi-signature of S1
founders or a jurisdictional $L_4$ claim, both of which terminate at
S1 or at jurisdictional authority (treated as an S1 surrogate).
Suppose for contradiction an infinite descending chain
$\identity_0 \succ \identity_1 \succ \identity_2 \succ \ldots$.
Each $\identity_i$ is recorded on chain with a block height $h_i$,
and the parent's registration precedes the child's, so
$h_0 > h_1 > h_2 > \ldots$. This is an infinite descending chain of
natural numbers --- impossible.
Hence finite-length termination at S1 (or an S1 surrogate) is
guaranteed.
\end{proof}

\Cref{thm:stratum-wellfounded} is the formal restatement of
``no accountability chain, no identity''
(\Cref{sec:formal-model}): every action by an S2/S4 identity traces in
finite steps to a human or jurisdictional accountability locus.

\subsection{Liability follows the chain}

\begin{corollary}[Liability attribution]\label{cor:liability}
For any action $\alpha$ taken by an identity $\identity$, liability
attaches to: (i)~$\identity$ itself if it carries a stake or bond;
(ii)~the parent identity in $\prec$ if $\identity$ is S2 or S4 and
unable to satisfy the liability; (iii)~recursively to the root S1
identity, by \Cref{thm:stratum-wellfounded}.
\end{corollary}

% =============================================================================
\section{LLM identity}\label{sec:llm}
% =============================================================================

LLMs are the test case that breaks naive substrate-equals-identity
reasoning: weights are bit-perfectly copyable.

\subsection{Substrate vs.\ instance}

\begin{definition}[LLM substrate]\label{def:llm-substrate}
An \emph{LLM substrate} is a tuple $(\mathbf{W}, \tau, A)$ where
$\mathbf{W}$ is a weights blob, $\tau$ is a tokenizer/architecture
specification, and $A$ is a training-run attestation. We write
$\Hash(\mathbf{W})$ for the canonical substrate hash.
\end{definition}

\begin{definition}[LLM instance]\label{def:llm-instance}
An \emph{LLM instance} is a tuple
$(\mathbf{W}_{\mathrm{ref}}, d, c)$ where $\mathbf{W}_{\mathrm{ref}}$
is a reference to a substrate hash, $d$ is a deployment identifier
(persistent across restarts), and $c$ is the conversational state.
\end{definition}

\subsection{Instance non-identity theorem}

\begin{theorem}[LLM instance non-identity]\label{thm:llm-non-identity}
Two LLM instances $\identity_1 = (\mathbf{W}, d_1, c_1)$ and
$\identity_2 = (\mathbf{W}, d_2, c_2)$ with identical substrate
reference but $d_1 \neq d_2$ are distinct identity loci.
\end{theorem}

\begin{proof}
By \Cref{def:identity-locus} an identity locus is determined by its
founding event $\founding$. By \Cref{def:llm-instance} the
deployment identifier $d$ is established at first inference call, hence
participates in $\founding$. Since $d_1 \neq d_2$, the founding events
$\founding_1 \neq \founding_2$, hence by
\Cref{thm:founding-uniqueness} the identity loci are distinct.
\end{proof}

\subsection{LLM tier model}

The four-tier permanence hierarchy of \Cref{sec:machine-storage} maps
to LLM identity:

\begin{table}[h]
\centering
\small
\begin{tabular}{@{}lp{10cm}@{}}
\toprule
\textsf{Tier} & \textsf{LLM storage location}\\
\midrule
0 & Weights hash; training-run Merkle root; architecture spec hash;
    TEE attestation (Intel TDX / AMD SEV-SNP / NVIDIA H100 confidential
    computing / Apple Private Cloud Compute) when applicable\\
1 & Deployment ID; LoRA adapter hash; system-prompt hash; tool/
    function-call registry\\
2 & API authentication; rate-limit/quota config; routing config;
    vector-DB or memory-store binding\\
3 & Conversation thread ID; context window; tool-call sequence; KV
    cache\\
\bottomrule
\end{tabular}
\caption{LLM identity by permanence tier (parallel to
\Cref{tab:tier-platforms}).}
\label{tab:llm-tiers}
\end{table}

\subsection{The substrate attestation pallet}

We introduce a candidate runtime pallet,
\texttt{pallet\_substrate\_attestation}, that records LLM substrate
provenance on chain. The extrinsic shape:

\begin{lstlisting}[language=rust]
pub fn attest_substrate(
    origin: OriginFor<T>,
    weights_hash: H256,
    training_run_merkle: H256,
    arch_spec_hash: H256,
    publisher_falcon_sig: BoundedVec<u8, ConstU32<2048>>,
    ipfs_cid_card: BoundedVec<u8, ConstU32<128>>,
) -> DispatchResult;
\end{lstlisting}

The pallet binds a weights hash to a training organization's
Falcon-512 public key. An LLM instance, on registration as an S4
agent, refers to a published substrate attestation. Verifiers can
confirm: the running agent claims the substrate that organization $O$
publicly attested to. The IPFS CID points to the model card and
training-run log; the on-chain record holds only hashes and the
publisher signature.

% =============================================================================
\section{Cryptographic substrate}\label{sec:crypto-substrate}
% =============================================================================

\subsection{The post-quantum imperative}

A classical signature is a time bomb. Suppose at time $t$ a signature
$\sigma = \Sign_{sk}(c)$ is produced under classical ECDSA. At time
$t' > t$, a CRQC becomes available. Any party with access to the
CRQC can compute a forgery $\sigma^\dagger$ on an arbitrary message
$c^\dagger$ that verifies under the same public key. From the
verifier's perspective at $t'$, $\sigma$ and $\sigma^\dagger$ are
indistinguishable: classical EUF-CMA security collapses retroactively.

Unlike encryption (where session keys can be rotated forward),
\emph{signatures bound to historical claims cannot be rotated}: once a
signature is part of the accretion chain, it is part of the record.

\subsection{Cryptographic survival theorem}

Let $\pqset$ denote the set of NIST-standardized post-quantum
signature schemes
($\pqset = \{\text{ML-DSA}, \text{SLH-DSA}, \text{FN-DSA}\}$ as of
2024--2026), and let $\Sigma_{\mathrm{cl}}$ denote classical schemes
(ECDSA, RSA-PSS, EdDSA).

\begin{theorem}[Cryptographic survival]\label{thm:cryptographic-survival}
Let $\identity$ be an identity locus with accretion chain $\accretion$
in which every $L_3$ claim uses some scheme $\Sigma \in \pqset$. Let
$T_{\mathrm{CRQC}}$ denote the (uncertain) time at which a
sufficiently-capable CRQC becomes available. Then for any verifier
inquiring at any time $t > T_{\mathrm{CRQC}}$, the
EUF-CMA-forgery probability of $\identity$'s accretion chain is
\[
\Pr[\text{forgery}] \leq \Negl(\lambda),
\]
where $\lambda$ is the post-quantum security parameter of $\Sigma$
(typically $\geq 128$).

Conversely, if any $L_3$ claim uses
$\Sigma' \in \Sigma_{\mathrm{cl}}$, then
\[
\Pr[\text{forgery}] \geq 1 - \Negl'(\lambda'),
\]
where $\Negl'$ characterizes the residual difficulty for the CRQC, and
the bound is dominated by the classical key length, not the
post-quantum security parameter.
\end{theorem}

\begin{proof}[Proof sketch]
The forward direction follows from the EUF-CMA security of
$\Sigma \in \pqset$ under the post-quantum adversarial model assumed
in NIST FIPS~203--205~\cite{nist-fips-203,nist-fips-204,nist-fips-205}.
The reverse direction follows from Shor's
algorithm~\cite{shor-1994}, which solves discrete-log in $\mathbb{Z}_p^*$
and elliptic-curve groups in polynomial quantum time, breaking ECDSA
and RSA-based signatures with overwhelming probability once the CRQC
operates above the corresponding key-size threshold.
\end{proof}

\begin{corollary}[Engineering implication]\label{cor:pq-must-now}
For any accretion chain intended to remain verifiable past
$T_{\mathrm{CRQC}}$, every $L_3$ claim must be signed under
$\Sigma \in \pqset$ \emph{at the time the claim is made}. Retrofitting
PQ signatures onto classical-signed historical claims is impossible:
the historical signatures already exist and are already forgeable.
\end{corollary}

\Cref{cor:pq-must-now} is the operational claim that \qh\ implements
SPHINCS+ (SLH-DSA, FIPS~205) and Falcon-512 (FN-DSA, FIPS draft~206) as
\emph{primary} signature schemes, not as opt-in upgrades. Every
identity-bearing transaction is PQ-signed at issuance.

\subsection{Signature size accommodation}

PQ signatures are larger than classical signatures. SPHINCS+-256f-simple
is 49{,}856~bytes; Ed25519 is 64~bytes. The \paraxiom\ fork of
\texttt{polkadot-sdk} raises \texttt{MAX\_SIGNATURE\_SIZE} to 50{,}000
and \texttt{parity-scale-codec}'s
\texttt{INITIAL\_PREALLOCATION} from 16~KiB to 256~KiB to accommodate
runtime metadata containing PQ signatures~\cite{paraxiom-mods-manifest}.
The trade-off (chain storage growth) is intentional: chain history
must persist; signature size is the cost of identity survival.

% =============================================================================
\section{Current digital identity landscape (May 2026)}\label{sec:landscape}
% =============================================================================

The formal model contacts deployed reality. We summarize the state of
digital identity at the time of writing, drawing on regulator filings,
standards publications, and operator disclosures. Detailed citations
are in the bibliography.

\subsection{Civic identity at scale}

\textbf{Estonia.} ID-card, Mobile-ID, Smart-ID. Classical cryptography
(RSA, ECDSA-P384). Cybernetica contracted in 2026 to author the
national PQC roadmap as part of the EU joint PQC plan. KSI-blockchain
audit log is hash-based and quantum-resilient; identity-binding
primitives are not.

\textbf{India --- Aadhaar.} 1.4~billion enrolled. Authentication
failure rate $\sim$6.5\% per attempt (Policy Circle, June~2025);
roughly 20.3~million failed attempts per month. 30~million ration
cards cancelled due to seeding mismatches. 815~million records
exposed in 2023; 29{,}000 biometric-cloning fraud incidents in 2024.
Architecturally a centralized database lookup, not a cryptographic
credential system. The canonical anti-pattern.

\textbf{EU --- eIDAS~2.0 / EUDI Wallet.} Regulation EU~2024/1183 in
force May~2024. Implementing acts December~2024. Hard deadline:
24~December~2026, every Member State must offer at least one EUDI
Wallet; mandatory acceptance by relying parties December~2027. Wire
protocols: OpenID4VP~1.0 Final, OpenID4VCI~1.0 Final, HAIP~1.0.
Credential formats: SD-JWT~VC and ISO~mdoc. PQC on a separate clock:
high-risk migration by 2030, full by 2035. Identity-wallet PKI
explicitly flagged as longest-lead-time transition.

\textbf{UK --- GOV.UK One~Login.} 11~million users. Becomes exclusive
access to central government services by end-2027. Currently lacks
certification under the UK Digital Identity \& Attributes Trust
Framework (DIATF) due to a supplier accreditation lapse --- the
national IDP operating uncertified against the country's own framework.

\textbf{US --- mDL.} No federal ID. State-level under AAMVA
Implementation Guidelines v1.4--1.5. As of January~2026, 21~states
plus Puerto Rico operate mDL programs accepted at TSA. ISO~mdoc with
COSE-signed claims. The only US system with cryptographic primitives
at meaningful scale.

\textbf{Canada.} Verified.me effectively defunct after SecureKey's
2022 sale. De-facto trust layer is DIACC's Pan-Canadian Trust
Framework (PCTF). No federal wallet equivalent to EUDI.

\textbf{Singapore, Brazil, Nigeria.} SingPass ($\sim$5M residents);
gov.br ($\sim$150M users with centralizing biometric service launched
2025); Nigeria NIN with biometric General Multipurpose Card from
January~2025. All architecturally centralized database-lookup; growing
scale, growing exclusion.

\textbf{Worldcoin / World~ID.} Banned, ordered to delete data, or under
active investigation in Colombia, Brazil, Thailand, Philippines, Hong
Kong, Spain, France, Portugal, Argentina, Kenya. Repeated
regulatory finding: financial incentives ($\$10$--$\$50$ in WLD per
iris scan) vitiate informed consent under GDPR-style biometric
regimes. Not a deployable architecture.

\subsection{Self-sovereign identity --- what shipped}

\begin{itemize}
\item W3C Verifiable Credentials Data Model 2.0 became a W3C
  Recommendation on 15~May~2025.
\item W3C DIDs Core 1.1 reached Candidate Recommendation Snapshot on
  5~March~2026; DID Resolution split into separate v0.3.
\item OpenID4VP~1.0 Final and OpenID4VCI~1.0 Final published
  July~2025; HAIP~1.0 Final December~2025. Selected as eIDAS~2.0
  ARF wire protocols.
\item Hyperledger Aries archived April~2025; components migrated to
  the OpenWallet Foundation.
\item \texttt{did:web} is the only DID method with material non-toy
  adoption (UN Transparency Protocol, Microsoft Entra Verified ID).
  \texttt{did:ion} lost momentum.
\end{itemize}

The deployed SSI surface is \texttt{did:web} plus OpenID4VP/VCI plus
SD-JWT~VC or ISO~mdoc. Everything else remains research-track.

\subsection{Blockchain-anchored identity}

ENS retains the only material adoption (1.7M+ domains; ENSv2 mainnet
February~2026). Sismo wound down late~2024. Spruce~ID pivoted to
government mDL integrations and now operates chain-agnostically.
Proof-of-personhood remains unsolved with no production winner.

\subsection{Hardware-rooted identity}

\begin{itemize}
\item WebAuthn / Passkeys / FIDO2: more than 1~billion users with at
  least one passkey activated; more than 15~billion accounts support
  passkeys. Apple/Google/Microsoft alignment held through 2025; FIDO
  CTAP~2.2 advanced cross-platform portability. HID/FIDO 2025
  enterprise survey reports 87\% of enterprises actively deploying
  passkeys.
\item DICE (Device Identifier Composition Engine, TCG specification):
  now the dominant root-of-trust pattern on constrained devices.
\item Apple Secure Enclave, Google Titan-M2, Microsoft Pluton: all
  baseline shipping on consumer devices.
\item Intel SGX: effectively abandoned for client-side use; deprecated
  on 11th-generation+ consumer CPUs. Server-side use shifted to Intel
  TDX and AMD SEV-SNP. ARM CCA Realms ships in Armv9.
\end{itemize}

\subsection{Post-quantum cryptography status}

\begin{itemize}
\item NIST FIPS~203 (ML-KEM), 204 (ML-DSA), 205 (SLH-DSA) published
  13~August~2024. FIPS~206 (FN-DSA / Falcon) in draft.
\item HQC announced March~2025 as second KEM (backup to ML-KEM).
\item US OMB PQC migration guidance, required within one year of NIST
  standards, was overdue as of May~2026: still in draft.
\item NSA CNSA~2.0: PQC required for National Security Systems by
  2030; classical algorithms prohibited by 2035.
\item EU joint PQC roadmap (June~2025): national strategies by
  end-2026; high-risk migration by 2030; full by 2035.
\item \textbf{No major civic identity system runs PQC signatures in
  production as of May~2026.}
\end{itemize}

\subsection{AI agent identity --- the open lane}

\begin{itemize}
\item Agentic AI Foundation (Linux Foundation) launched
  December~2025. Anthropic contributed MCP; OpenAI contributed
  AGENTS.md; Block contributed Goose.
\item MCP authorization standardized OAuth 2.1 + PKCE in
  March~2025; the November~2025 release added a \emph{server
  identity} primitive.
\item Google A2A protocol (April~2025) uses workload-identity
  federation: short-lived cryptographically verifiable credentials.
\item IETF Web Bot Auth draft (Cloudflare-backed) signs HTTP requests
  with private keys plus a published directory.
\item GNAP (RFC~9635, October~2024): most plausible OAuth successor
  for delegated agent authorization; minimal vendor uptake.
\item Model substrate attestation: \emph{no agreed standard}. NVIDIA
  H100/B200, AMD SEV-SNP, Intel TDX supply hardware; the protocol
  layer is open.
\end{itemize}

The state of AI agent identity in 2026 is comparable to the state of
SSI in 2018: many overlapping proposals, no production winner, vendors
shipping ahead of standards.

\subsection{Positioning fit for \paraxiom}

\begin{table}[h]
\centering
\small
\begin{tabular}{@{}lp{5cm}p{5cm}c@{}}
\toprule
\textsf{Layer} & \textsf{Industry default 2026} & \textsf{\paraxiom\ choice} & \textsf{Lane}\\
\midrule
Wallet protocol & OID4VP/VCI + HAIP & Implement OID4VP-compatible verifier & Track\\
Credential format & SD-JWT~VC + ISO mdoc & Support both for export & Track\\
DID method & \texttt{did:web} for interop & \texttt{did:web} + QH-anchored method & Track + extend\\
Signature crypto & RSA / ECDSA (classical) & SPHINCS+ / Falcon-512 (PQ) & \textbf{Lead}\\
Hardware root & Apple SE / Titan-M2 / Pluton & Same + DICE for edge & Track\\
Agent identity & MCP + OAuth 2.1 (early) & MCP + Falcon-signed + on-chain & \textbf{Lead}\\
Model attestation & None (open) & \texttt{pallet\_substrate\_attestation} & \textbf{Lead}\\
Trust framework & EU eIDAS / DIACC / state & Chain-native, bridge to civic & Sovereign\\
\bottomrule
\end{tabular}
\caption{\paraxiom\ positioning against the 2026 landscape. Three
uncontested lanes: post-quantum civic-grade identity, agent
cryptographic identity, model substrate attestation.}
\label{tab:positioning}
\end{table}

% =============================================================================
\section{The \paraxiom\ architecture}\label{sec:paraxiom-architecture}
% =============================================================================

We now describe the engineering realization.

\subsection{Three-tier system}

The \paraxiom\ identity stack is decomposed into three layers of
storage with three distinct trust and performance profiles:

\begin{table}[h]
\centering
\begin{tabular}{@{}llll@{}}
\toprule
\textsf{Layer} & \textsf{Holds} & \textsf{Speed} & \textsf{Mutability}\\
\midrule
\qh\ chain & Identity records, hashes, CIDs, commitments, signatures & Slow & Append-only\\
IPFS & Encrypted payloads, KYC bundles, vault ciphertext & Medium & Immutable (by CID)\\
Postgres + Redis & Denormalized indexes, query caches & Fast & Rebuildable mirror\\
\bottomrule
\end{tabular}
\caption{The three storage layers of the \paraxiom\ identity stack.}
\label{tab:three-tier}
\end{table}

The chain is the source of truth for cryptographic claims. IPFS is
the canonical bulk-payload store, content-addressed by CID. The
Postgres+Redis mirror is rebuildable from chain plus IPFS without
information loss (modulo unrecoverable plaintexts, which are protected
by deletion of the master key per \Cref{thm:crypto-shredding}).

\subsection{Identity strata}

The five strata of \Cref{def:stratification} are realized on chain
through the existing
\texttt{identity-registry} pallet (extended) and a small number of
companion pallets:

\begin{lstlisting}[language=rust]
pub struct Identity<AccountId, BlockNumber> {
    pub id: u64,
    pub owner: AccountId,
    pub public_key: BoundedVec<u8, ConstU32<2048>>,
    pub role: IdentityRole,
    pub status: IdentityStatus,
    pub reputation_bps: u32,
    pub created_at_block: BlockNumber,
    pub created_by_tx_hash: H256,
    pub key_rotation_history: BoundedVec<KeyRotation, ConstU32<32>>,

    /// Stratum tag (S1, S2, S4, S5).
    pub stratum: Stratum,

    /// Parent identity for S2 and S4 (required by Theorem 6.2).
    pub parent_identity: Option<u64>,

    /// Tenant context for S2.
    pub tenant_id: Option<BoundedVec<u8, ConstU32<64>>>,

    /// Custody mode (SelfSovereign / NcCustodial / TenantCustodial).
    pub custody: CustodyMode,

    /// KYC bundle hash for S1 (Falcon-signed by KYC issuer).
    pub kyc_bundle_hash: Option<[u8; 32]>,

    /// Instance anchor for S4 agents (per Corollary 5.1).
    pub instance_anchor: Option<[u8; 32]>,
}
\end{lstlisting}

\subsection{Companion pallets}

\begin{description}
\item[\texttt{pallet-onboarding-commit}] (extended): a generalized
  commit-reveal flow with bond and slashable dispute window. Used to
  bootstrap new S1 identities, agent S4 identities, and validators.
\item[\texttt{axiom-attestation}] (extended): on-chain Falcon-signed
  attestations with IPFS CIDs for full payloads. Extended with a
  \texttt{VerifierKind} enum to dispatch question-generation,
  PR-code-review, vault-commitment, persona-attestation, and existing
  Axiom-task variants.
\item[\texttt{pallet-credential-vault}] (new): on-chain commitments to
  off-chain envelope-encrypted credentials. Stores
  \texttt{(ciphertext\_cid, ciphertext\_hash, iv\_hint, service, kind,
  created\_at, revoked\_at)}.
\item[\texttt{pallet-reputation-tokens}] (new): soulbound earned
  reputation tokens with minting policies indexed by behavior. Distinct
  from the governance-adjustable \texttt{reputation\_bps} of the
  identity registry.
\item[\texttt{pallet-substrate-attestation}] (new): per
  \Cref{sec:llm}, publishes LLM weights provenance on chain.
\end{description}

\subsection{Credential vault flow}

\begin{figure}[h]
\centering
\begin{tikzpicture}[every node/.style={font=\small},
                    box/.style={draw=ncrule,fill=nclight,
                                rounded corners=2pt,
                                minimum width=4cm,
                                minimum height=0.7cm,
                                inner sep=4pt},
                    arrow/.style={->,>=Stealth,thick,ncaccent}]
  \node[box] (a) at (0,4)  {Plaintext credential};
  \node[box] (b) at (0,3)  {Envelope: $iv \| tag \| ct$};
  \node[box] (c) at (0,2)  {IPFS pin $\rightarrow$ CID};
  \node[box] (d) at (0,1)  {On-chain commitment record};
  \draw[arrow] (a) -- node[right]{AES-256-GCM with user master key} (b);
  \draw[arrow] (b) -- node[right]{$\Hash$ and pin} (c);
  \draw[arrow] (c) -- node[right]{Extrinsic with CID + hash} (d);
\end{tikzpicture}
\caption{Credential vault encryption flow. Plaintext never leaves the
client (\emph{v0.2+}) or server boundary (\emph{v0.1}); IPFS holds
ciphertext only; chain holds commitment.}
\label{fig:vault-flow}
\end{figure}

\subsection{Crypto-shredding completeness}

A key correctness property of the credential vault: revocation through
master-key deletion must render all encrypted credentials inaccessible,
\emph{even if the ciphertext remains pinned to IPFS}.

\begin{theorem}[Crypto-shredding completeness]\label{thm:crypto-shredding}
Let $k$ be a master key over a strongly-secure AEAD scheme (AES-256-GCM
under the standard model). Let $\{C_1, \ldots, C_n\}$ be a set of
ciphertexts $C_i = \Enc_k(m_i, iv_i)$ for messages $m_i$ and unique
nonces $iv_i$. Suppose $k$ is securely erased from all storage
controlled by the holder. Then, for any computationally-bounded
adversary $\mathcal{A}$ with access to all $C_i$ but not to $k$,
\[
\Pr\bigl[\mathcal{A}(C_1, \ldots, C_n) = m_i \bigr] \leq \frac{1}{|\mathcal{M}|} + \Negl(\lambda),
\]
where $\mathcal{M}$ is the plaintext domain and $\lambda$ is the AEAD
security parameter.
\end{theorem}

\begin{proof}
By the IND-CPA security of AES-256-GCM under the standard model.
$\Negl(\lambda)$ is the adversary's IND-CPA advantage at security
parameter $\lambda = 256$. Crucially, AES-256-GCM's security guarantee
binds the security to key secrecy, not to ciphertext exclusivity:
public ciphertext does not weaken the bound. Erasing $k$ thus removes
all information-theoretic distinguishability from the chosen-message
oracle.
\end{proof}

\begin{corollary}[CCPA crypto-shredding]\label{cor:ccpa}
Compliance with right-to-erasure obligations (CCPA \S1798.105, GDPR
Art.~17) can be discharged by deleting the user's master key from all
holder-controlled stores, without requiring unpinning of IPFS-resident
ciphertexts.
\end{corollary}

\subsection{Sync mechanism}

The Postgres+Redis mirror in the gov-app subscribes to chain events
via Substrate's WebSocket RPC. Events from \texttt{identity-registry},
\texttt{axiom-attestation}, \texttt{pallet-reputation-tokens}, and
\texttt{pallet-credential-vault} are decoded and applied to the
relational schema; Redis caches are invalidated on the keyed indexes.

\subsection{Verifier at action time}

\begin{theorem}[Verification-time accountability]\label{thm:verify-action-time}
Let $\identity$ be an identity attempting an action $\alpha$ at time
$t$. The \texttt{verify\_chain}$(\identity, \alpha, t)$ helper returns
$\top$ iff (i)~$\identity$ has an active $L_3$ key at $t$, (ii)~$\identity$'s
$\prec$-ancestry terminates in finite steps at an S1 with a non-empty
$L_4$ claim or in a stake-bonded entity, and (iii)~$\identity$'s
$L_5$ reputation satisfies the action's minimum threshold $\theta_\alpha$.
\end{theorem}

\begin{proof}
Conditions (i) and (iii) are checked by direct query against
\texttt{identity-registry} and \texttt{pallet-reputation-tokens}.
Condition~(ii) terminates in finite steps by
\Cref{thm:stratum-wellfounded}.
\end{proof}

% =============================================================================
\section{Future scenarios}\label{sec:future}
% =============================================================================

We project three scenarios for the trajectory of identity past 2030.

\subsection{Dissolution}

Identity collapses if (a)~AI agents proliferate faster than
accountability chains can be maintained; (b)~CRQC arrives before
PQC migration completes for civic-identity infrastructure;
(c)~states lose capacity to issue civic identity; (d)~privacy norms
preclude reputation tracking. Non-zero probability; the canonical
reason this work is timely.

\subsection{Fragmentation}

Identity persists but jurisdictionalizes: each major bloc operates
its own infrastructure with limited interoperability. This is the
trajectory the world is presently on. \paraxiom's chain-native,
state-agnostic substrate functions as a bridge across fragmented
zones, attesting to jurisdictional claims when needed via
\texttt{kyc\_bundle\_hash} but not depending on any single
jurisdiction's continued existence.

\subsection{Augmentation}

Identity extends: humans, agents, swarms, organizations, protocols
all carry first-class cryptographic identity with explicit
accountability chains and portable reputation. The five-stratum design
of \Cref{def:stratification} is a first draft of an augmented identity
ontology; further strata (swarm, protocol, hybrid) extend the schema.

The outcome among these three depends not on inevitability but on
what is built. Dissolution occurs by default if no one builds the
infrastructure; fragmentation if states build it without
interoperability; augmentation if deliberate choice produces
post-quantum, decentralized, accountability-preserving substrates.

% =============================================================================
\section{Implementation roadmap}\label{sec:roadmap}
% =============================================================================

We sequence the engineering work.

\begin{enumerate}
\item \textbf{Identity-registry extensions} ($\sim$3 days).
  Non-breaking; ships first. Adds stratum, parent\_identity, tenant\_id,
  custody, kyc\_bundle\_hash, instance\_anchor fields.
\item \textbf{pallet-onboarding-commit generalization} ($\sim$2 days).
  Per-role bonds; supports S1/S4/validator commit-reveal.
\item \textbf{axiom-attestation VerifierKind extension} ($\sim$2 days).
  Adds QuestionGen, PrCodeReview, VaultCommitment, PersonaAttestation
  variants alongside the existing AxiomTask.
\item \textbf{pallet-credential-vault} ($\sim$3 days). On-chain
  commitments to IPFS-resident encrypted credentials. Realizes
  \Cref{thm:crypto-shredding}.
\item \textbf{pallet-reputation-tokens} ($\sim$4 days). Earned,
  soulbound, threshold-minted tokens.
\item \textbf{pallet-substrate-attestation} ($\sim$3 days). Per
  \Cref{sec:llm}, publishes LLM substrate provenance.
\item \textbf{Gov-app event subscriber} ($\sim$5 days). Subxt or
  Polkadot.js client; consumes chain events; populates Postgres mirror.
\item \textbf{Spec version bump + on-chain migration} ($\sim$3 days).
  Migrates existing identities to the extended schema.
\end{enumerate}

Total estimated effort: 3--4 weeks of focused work, parallelizable
with gov-app v0.2 work that depends on the migrated schema.

% =============================================================================
\section{Conclusion}\label{sec:conclusion}
% =============================================================================

Identity is engineering. The accountability chain does not extend
itself; the post-quantum signatures do not sign themselves; the
strata do not enumerate themselves. The substrates we now share with
machines and AI agents will, by the early 2030s, demand identity
infrastructure that survives a quantum-cryptographic transition, that
scales without requiring social acquaintance, that holds open the
possibility of new kinds of accountability loci we cannot fully
enumerate today.

We have stated identity formally as a layered accountability chain
above a founding individuation event, proven four invariants
(founding uniqueness, tier monotonicity, cryptographic survival,
crypto-shredding completeness), surveyed the 2026 deployed landscape
and located \paraxiom's three uncontested lanes (PQ civic-grade,
agent cryptographic identity, model substrate attestation), and
specified the on-chain plus IPFS plus mirror architecture that
realizes the model.

The question with which this treatise was commissioned --- ``will we
have identity in the future?'' --- admits a precise answer. Yes, if
built. The pallets specified in \Cref{sec:paraxiom-architecture} are
the construction.

% =============================================================================
\appendix

\section{Detailed proofs}\label{app:proofs}

\subsection{Founding uniqueness (full)}

\begin{proof}[Full proof of \Cref{thm:founding-uniqueness}]
We argue from \Cref{def:identity-locus}, \Cref{def:founding-event},
and \Cref{ax:founding-precedence}.

Let $\identity = (\substrate, \founding, \accretion, \relations)$ be
an identity locus and suppose for contradiction that there exist two
founding-event candidates $\founding_1 = (s, t_0^{(1)}, \pi_1)$ and
$\founding_2 = (s, t_0^{(2)}, \pi_2)$, both consistent with
$\identity$ in the sense that every claim in $\accretion$ post-dates
the candidate's timestamp.

Case 1: $t_0^{(1)} < t_0^{(2)}$. Then there exist points in time
$t \in (t_0^{(1)}, t_0^{(2)})$ at which $\substrate$ was already
individuated relative to $\founding_1$ but, by hypothesis, not yet
individuated relative to $\founding_2$. Any claim made in this
interval --- and \Cref{ax:founding-precedence} permits such claims to
exist in $\accretion$ relative to $\founding_1$ --- would violate
\Cref{ax:founding-precedence} relative to $\founding_2$. Either
$\founding_2$ does not consistently anchor $\identity$, or no claims
exist in the interval (in which case $t_0^{(2)}$ can be taken as
$t_0^{(1)}$ without loss of information). In both sub-cases we
collapse to a single founding event.

Case 2: $t_0^{(1)} = t_0^{(2)}$. Then both events name the same
moment of individuation of the same substrate; up to the witness set,
they coincide. The witness sets may be merged: $\pi = \pi_1 \cup
\pi_2$ produces the most informative single $\founding$, and the
two-event presentation reduces to a single-event presentation.

In both cases the two candidates collapse to one, contradicting the
assumption of distinctness.
\end{proof}

\subsection{Stratum well-foundedness (full)}

\begin{proof}[Full proof of \Cref{thm:stratum-wellfounded}]
We rely on three facts about the chain registration protocol:
\begin{enumerate}
\item Every S2 or S4 registration extrinsic requires a non-null
  \texttt{parent\_identity} field; the chain validates that the named
  parent exists prior to applying the registration.
\item S5 registration requires either a multi-signature attestation
  of S1 founders or a jurisdictional $L_4$ claim. Both terminate
  finitely at S1 or at a jurisdictional authority treated as an S1
  surrogate.
\item Every identity in \texttt{identity-registry} carries a
  \texttt{created\_at\_block} field, strictly less than the block in
  which any of its children are registered.
\end{enumerate}

Suppose for contradiction an infinite descending chain
$\identity_0 \succ \identity_1 \succ \identity_2 \succ \ldots$. By
fact~3, $\mathrm{block}(\identity_0) > \mathrm{block}(\identity_1) >
\mathrm{block}(\identity_2) > \ldots$, an infinite strictly
decreasing sequence in $\mathbb{N}$, which is impossible.

Thus every descending chain terminates in finite steps. By facts~1
and~2, the terminus is either an S1 or an S1 surrogate. In either
case, finite-step termination at an S1-equivalent is guaranteed.
\end{proof}

\subsection{Cryptographic survival (full)}

\begin{proof}[Full proof of \Cref{thm:cryptographic-survival}]
\textbf{Forward direction.} For $\Sigma \in \pqset$ each scheme has
been demonstrated EUF-CMA-secure against any polynomial-time quantum
adversary under the standard model, per the security reductions
underlying NIST FIPS~203--205. Let $\mathcal{Q}$ be such a quantum
adversary with running time bounded by $\mathrm{poly}(\lambda)$. The
EUF-CMA security definition states:
\[
\Pr[\mathrm{EUF\text{-}CMA}_{\mathcal{Q}}(\Sigma, \lambda) = 1] \leq \Negl(\lambda).
\]
This implies that the probability of any forgery against any single
claim $a_i$ in $\accretion$ is at most $\Negl(\lambda)$. By a union
bound over the $n$ claims, the probability of any forgery anywhere
in $\accretion$ is at most $n \cdot \Negl(\lambda)$, which remains
negligible for any polynomial-size $\accretion$.

\textbf{Reverse direction.} For $\Sigma' \in \Sigma_{\mathrm{cl}}$,
Shor's algorithm~\cite{shor-1994} solves the discrete logarithm
problem in $\mathbb{Z}_p^*$ and on elliptic curves over $\mathbb{F}_p$
in time polynomial in $\log p$. Given a CRQC of size sufficient to
implement Shor at the relevant key size, the adversary can recover
$sk$ from $pk$ with probability $1 - \Negl'(\lambda')$, after which
forgery is trivial via $\Sign_{sk}$. The probability bound on forgery
is hence at least $1 - \Negl'(\lambda')$, where $\Negl'(\lambda')$ is
characterized by the CRQC's residual error, dominated by the
classical key length $\lambda'$.

\end{proof}

\subsection{Crypto-shredding completeness (full)}

\begin{proof}[Full proof of \Cref{thm:crypto-shredding}]
We reduce to the IND-CPA security of AES-256-GCM. Suppose for
contradiction that $\mathcal{A}$ achieves
\[
\Pr[\mathcal{A}(C_1, \ldots, C_n) = m_i] > \frac{1}{|\mathcal{M}|} + \epsilon
\]
for non-negligible $\epsilon$, for some target index $i$. Construct
an IND-CPA adversary $\mathcal{B}$ against AES-256-GCM as follows:
$\mathcal{B}$ receives an IND-CPA challenge oracle, queries it on
all $\{m_1, \ldots, m_n\}$ to obtain ciphertexts identically
distributed to those $\mathcal{A}$ would receive, then runs
$\mathcal{A}$ on these ciphertexts. The output of $\mathcal{B}$ is
$\mathcal{A}$'s output.

If $\mathcal{A}$ recovers $m_i$ with advantage $\epsilon$ over the
$1/|\mathcal{M}|$ baseline, then $\mathcal{B}$ distinguishes the
IND-CPA challenge with the same advantage. By IND-CPA security of
AES-256-GCM, this advantage is bounded by $\Negl(\lambda)$ where
$\lambda = 256$ is the AES key length. Hence $\epsilon \leq
\Negl(\lambda)$, contradicting the supposed non-negligibility.
\end{proof}

\section{Pallet API specifications}\label{app:pallet-api}

\subsection{\texttt{identity-registry} (extended)}

\begin{lstlisting}[language=rust]
pub enum Stratum { S1Person, S2Scoped, S4AgentInstance, S5CrossTenant }

pub enum CustodyMode { SelfSovereign, NcCustodial, TenantCustodial }

pub enum IdentityRole {
    Shield, Agent, OracleReporter, Validator,
    Citizen, Decideur, Fournisseur, Service,
}

pub enum IdentityStatus {
    Active, Suspended, Revoked, Deceased, Decommissioned, Dissolved
}

/// Register an S1 identity, optionally with KYC bundle.
pub fn register_identity_s1(
    origin: OriginFor<T>,
    public_key: BoundedVec<u8, ConstU32<2048>>,
    kyc_bundle_hash: Option<[u8; 32]>,
    custody: CustodyMode,
) -> DispatchResult;

/// Delegate to a scoped S2 identity from a caller S1.
pub fn delegate_to_scope(
    origin: OriginFor<T>,
    tenant_id: BoundedVec<u8, ConstU32<64>>,
    role: IdentityRole,
) -> DispatchResult;

/// Register an S4 agent instance under a parent identity.
pub fn register_agent_instance(
    origin: OriginFor<T>,
    parent_identity: u64,
    agent_kind: BoundedVec<u8, ConstU32<64>>,
    public_key: BoundedVec<u8, ConstU32<2048>>,
    instance_anchor: [u8; 32],
) -> DispatchResult;

/// Verify the accountability chain for an action.
pub fn verify_chain(
    identity_id: u64,
    action_kind: ActionKind,
    block: BlockNumberFor<T>,
) -> Result<ChainVerdict, DispatchError>;
\end{lstlisting}

\subsection{\texttt{pallet-credential-vault}}

\begin{lstlisting}[language=rust]
pub struct Commitment<BlockNumber> {
    pub holder_identity: u64,
    pub service_name: BoundedVec<u8, ConstU32<64>>,
    pub credential_kind: CredentialKind,
    pub ciphertext_cid: BoundedVec<u8, ConstU32<128>>,
    pub ciphertext_hash: [u8; 32],
    pub iv_hint: [u8; 12],
    pub created_at: BlockNumber,
    pub revoked_at: Option<BlockNumber>,
}

pub fn commit(
    origin: OriginFor<T>,
    service_name: BoundedVec<u8, ConstU32<64>>,
    credential_kind: CredentialKind,
    ciphertext_cid: BoundedVec<u8, ConstU32<128>>,
    ciphertext_hash: [u8; 32],
    iv_hint: [u8; 12],
) -> DispatchResult;

pub fn revoke(
    origin: OriginFor<T>,
    commitment_id: u64,
) -> DispatchResult;

pub fn attest_use(
    origin: OriginFor<T>,
    commitment_id: u64,
    caller_context: BoundedVec<u8, ConstU32<128>>,
) -> DispatchResult;
\end{lstlisting}

\subsection{\texttt{pallet-substrate-attestation}}

\begin{lstlisting}[language=rust]
pub struct SubstrateAttestation<BlockNumber> {
    pub weights_hash: H256,
    pub training_run_merkle: H256,
    pub arch_spec_hash: H256,
    pub publisher: u64,                // identity-registry ID
    pub publisher_sig: BoundedVec<u8, ConstU32<2048>>,
    pub ipfs_cid_card: BoundedVec<u8, ConstU32<128>>,
    pub published_at: BlockNumber,
    pub revoked_at: Option<BlockNumber>,
}

pub fn attest_substrate(
    origin: OriginFor<T>,
    weights_hash: H256,
    training_run_merkle: H256,
    arch_spec_hash: H256,
    publisher_falcon_sig: BoundedVec<u8, ConstU32<2048>>,
    ipfs_cid_card: BoundedVec<u8, ConstU32<128>>,
) -> DispatchResult;
\end{lstlisting}

\section{Sources}\label{app:sources}

\begin{thebibliography}{99}

\bibitem{nist-fips-203} National Institute of Standards and Technology.
\emph{FIPS~203: Module-Lattice-Based Key-Encapsulation Mechanism
Standard.} U.S.\ Department of Commerce, 13~August~2024.

\bibitem{nist-fips-204} National Institute of Standards and Technology.
\emph{FIPS~204: Module-Lattice-Based Digital Signature Standard.}
U.S.\ Department of Commerce, 13~August~2024.

\bibitem{nist-fips-205} National Institute of Standards and Technology.
\emph{FIPS~205: Stateless Hash-Based Digital Signature Standard.}
U.S.\ Department of Commerce, 13~August~2024.

\bibitem{shor-1994} P.~W. Shor. \emph{Algorithms for quantum
computation: discrete logarithms and factoring.} Proceedings of the
35th Annual Symposium on Foundations of Computer Science (FOCS),
1994.

\bibitem{eidas-2024-1183} European Parliament and Council. \emph{Regulation
(EU)~2024/1183 amending Regulation (EU) No~910/2014 as regards
establishing the European Digital Identity Framework (eIDAS~2.0).}
Official Journal of the European Union, May~2024.

\bibitem{mcp-spec-2025-11} Anthropic. \emph{Model Context Protocol
Specification, Release 2025-11-25.} \url{https://modelcontextprotocol.io/}

\bibitem{google-a2a} Google. \emph{Agent-to-Agent (A2A) Protocol.}
Published April~2025.

\bibitem{paraxiom-mods-manifest} Paraxiom Technologies Inc.\
\emph{PARAXIOM-MODIFICATIONS.md: canonical fork-change manifest for
polkadot-sdk and parity-scale-codec.} \qh\ repository, 2026.

\bibitem{paraxiom-identity-registry} Paraxiom Technologies Inc.\
\emph{pallet-identity-registry source code.} \qh\ pallets directory,
2026.

\bibitem{aamva-mdl-15} American Association of Motor Vehicle
Administrators. \emph{Mobile Driver License Implementation Guidelines
v1.5.} May~2026.

\bibitem{w3c-vcdm-2} World Wide Web Consortium. \emph{Verifiable
Credentials Data Model 2.0 (W3C Recommendation).} 15~May~2025.

\bibitem{w3c-did-1-1} World Wide Web Consortium. \emph{Decentralized
Identifiers (DIDs) v1.1 (W3C Candidate Recommendation Snapshot).}
5~March~2026.

\bibitem{openid-vp-1} OpenID Foundation. \emph{OpenID for Verifiable
Presentations 1.0 (Final).} July~2025.

\bibitem{openid-vci-1} OpenID Foundation. \emph{OpenID for Verifiable
Credential Issuance 1.0 (Final).} July~2025.

\bibitem{nsa-cnsa-2} National Security Agency. \emph{Commercial
National Security Algorithm Suite 2.0.} U.S.\ Department of Defense,
September~2022 (updated 2025).

\bibitem{enisa-pqc} European Union Agency for Cybersecurity. \emph{Post-
Quantum Cryptography: Current State and Quantum Mitigation.} ENISA,
2025.

\bibitem{tcg-dice} Trusted Computing Group. \emph{Device Identifier
Composition Engine (DICE) Architectures.} TCG Specification,
2024 revision.

\bibitem{policy-circle-aadhaar} Policy Circle. \emph{Aadhaar authentication
failures persist after a decade.} June~2025.

\bibitem{linux-foundation-aaif} Linux Foundation. \emph{Agentic AI
Foundation launch announcement.} December~2025.

\end{thebibliography}

\bigskip

\noindent\emph{Document version}: 1.0, May~15,~2026.
\emph{Companion materials}:
\texttt{IDENTITY-ARCHITECTURE.md} (narrative version);
\texttt{IDENTITY-PALLETS-DESIGN.md} (engineering specification);
\texttt{GOVERNANCE-ROADMAP.md} \S{}10 (product roadmap).
\emph{Repository}: \url{https://github.com/Paraxiom/quantumharmony} and
companion documents in the NeoCarbone lovable-bridge repository.

\end{document}
